\documentclass{beamer}
\usepackage[utf8]{inputenc}

\usetheme{Madrid}
\usecolortheme{default}
\usepackage{amsmath,amssymb,amsfonts,amsthm}
\usepackage{txfonts}
\usepackage{tkz-euclide}
\usepackage{listings}
\usepackage{adjustbox}
\usepackage{array}
\usepackage{tabularx}
\usepackage{gvv}
\usepackage{lmodern}
\usepackage{circuitikz}
\usepackage{tikz}
\usepackage[utf8]{inputenc}

\usetheme{Madrid}
\usecolortheme{default}
\usepackage{amsmath,amssymb,amsfonts,amsthm}
\usepackage{txfonts}
\usepackage{tkz-euclide}
\usepackage{listings}
\usepackage{adjustbox}
\usepackage{array}
\usepackage{tabularx}
\usepackage{gvv}
\usepackage{lmodern}
\usepackage{circuitikz}
\usepackage{tikz}
\usepackage{graphicx}

\setbeamertemplate{page number in head/foot}[totalframenumber]

\usepackage{tcolorbox}
\tcbuselibrary{minted,breakable,xparse,skins}



\definecolor{bg}{gray}{0.95}
\DeclareTCBListing{mintedbox}{O{}m!O{}}{%
  breakable=true,
  listing engine=minted,
  listing only,
  minted language=#2,
  minted style=default,
  minted options={%
    linenos,
    gobble=0,
    breaklines=true,
    breakafter=,,
    fontsize=\small,
    numbersep=8pt,
    #1},
  boxsep=0pt,
  left skip=0pt,
  right skip=0pt,
  left=25pt,
  right=0pt,
  top=3pt,
  bottom=3pt,
  arc=5pt,
  leftrule=0pt,
  rightrule=0pt,
  bottomrule=2pt,

  colback=bg,
  colframe=orange!70,
  enhanced,
  overlay={%
    \begin{tcbclipinterior}
    \fill[orange!20!white] (frame.south west) rectangle ([xshift=20pt]frame.north west);
    \end{tcbclipinterior}},
  #3,
}
\lstset{
    language=C,
    basicstyle=\ttfamily\small,
    keywordstyle=\color{blue},
    stringstyle=\color{orange},
    commentstyle=\color{green!60!black},
    numbers=left,
    numberstyle=\tiny\color{gray},
    breaklines=true,
    showstringspaces=false,
}
%------------------------------------------------------------
%This block of code defines the information to appear in the
%Title page
\title %optional
{5.13.40}
\date{October  2025}
%\subtitle{A short story}

\author % (optional)
{Bhoomika L - EE25BTECH11014}

\begin{document}


\frame{\titlepage}
\begin{frame}{Question} Let $a,\lambda,\mu \in \mathbb{R}$. Consider the system of linear equations\\
\begin{center}
$ax+2y=\lambda$\\
$3x-2y=\mu$
\end{center}
Which of the following statement$\brak{s}$ is $\brak{are}$ correct?
\begin{enumerate}
\item If $a=-3,$ then the system has infinitely many solutions for all value of $\lambda$ and $\mu$.
 \item If $a \neq -3$, then the system has unique solution for all values of $\lambda$ and $\mu.$
 \item  If $\lambda + \mu = 0$, then the system has infinitely many solutions for $a = -3$.
\item  If $\lambda + \mu \neq 0$, then the system has no solution for $a = -3$.
\end{enumerate}
\end{frame}

\begin{frame}{Solution}
By the system of equations,we get the following augmented matrix:\\
By row reduction method,
\begin{align} 
\myvec{\begin{array}{cc|c}
a & 2 & \lambda \\
3 & -2 & \mu
\end{array}} \xrightarrow{R_2\longrightarrow R_2-3\frac{R_1}{a}}\myvec{\begin{array}{cc|c}
a & 2 & \lambda \\
0 & \frac{-2a-6}{a} & \mu-\frac{3\lambda}{a}
\end{array}} \\
Now \ if \ \frac{-2a-6}{a}=0 \implies a=-3,\\
when \ a=-3,\mu-3\frac{\lambda}{a}=\mu+\lambda
\end{align}

\textbf{Case-1:}If $\lambda+\mu=0$ and the system has \textbf{infinitely many solutions}.\\
\textbf{Case-2:}If $\lambda+\mu\neq0$, then the system has \textbf{no solution}\\

\end{frame}

\begin{frame}{Solution}
and if $a\neq-3$, then the first row is non-zero and matrix is invertible. Hence system has a \textbf{unique solution} for all values of $\lambda$ and $\mu$. \\

And solution for $a\neq-3$ is:
\begin{align}
   x\brak{a+3}=\lambda+\mu \implies x=\frac{\lambda+\mu}{a+3}\\
   3x-2y=\mu \implies y=\frac{3\lambda-a\mu}{2a+6}
 \end{align}

Therefore (2),(3) and (4) are the correct statements.    
\end{frame}


 \begin{frame}[fragile]
\frametitle{C Code }
\begin{lstlisting}
#include <stdio.h>

void solveSystem(double a, double lambda, double mu, double sol[2]) {
    double D = -2 * (a + 3);

    if (D != 0) {
        sol[0] = (-2 * lambda - 2 * mu) / D;  // x
        sol[1] = (3 * lambda - a * mu) / D;   // y
        printf("Unique solution exists:\n");
        printf("x = %.3f, y = %.3f\n", sol[0], sol[1]);
    } else {
        // When a = -3
        if (lambda + mu == 0)
            printf("Infinitely many solutions (system consistent)\n");
        else
            printf("No solution (system inconsistent)\n");
    }
}
\end{lstlisting}
\end{frame}

 \begin{frame}[fragile]
\frametitle{Python Code }
\begin{lstlisting}
import ctypes

# Load shared library
lib = ctypes.CDLL("./code.so")

# Define argument and return types
lib.solveSystem.argtypes = [
    ctypes.c_double, ctypes.c_double, ctypes.c_double,
    ctypes.POINTER(ctypes.c_double)
]
\end{lstlisting}
\end{frame}

 \begin{frame}[fragile]
\frametitle{Python Code }
\begin{lstlisting}
# Prepare result array
sol = (ctypes.c_double * 2)()

# ---- Fixed values ----
a = -3
lam = 2
mu = -2

# ---- Call C function ----
lib.solveSystem(a, lam, mu, sol)

# ---- Display result (if unique) ----
if a != -3:
    print(f"x = {sol[0]:.3f}, y = {sol[1]:.3f}")

\end{lstlisting}
\end{frame}

\end{document}